\documentclass[compress, aspectratio=169]{beamer}

\usepackage{presentationtemplate}
\usepackage[fontsize=\small,linenosfontsize=\small]{mylisting}
\usepackage[fontsize=\small]{terminal}

\title{GNU Make: от shell-скрипта к инкрементальной сборке}

\setbeamertemplate{footnote}{%
  \small\insertfootnotetext\par
}

% Beamer ставит перед каждой сноской правило высотой
% \footnotesep (по умолчанию 10pt). Уменьшаем отступ между
% сносками: в кадре 33 таблица и три сноски переполняли
% высоту кадра на 2,8pt при стандартном значении.
\setlength{\footnotesep}{3pt}

% Базовый URL руководства GNU Make.
% Тело разбито на строки (каждая < 100 символов), символы %
% гасят переводы строк, чтобы в URL не попало пробельное
% пространство. Внутри \url{...} макрос расширяется в полный
% адрес (проверено: и вывод, и цель ссылки корректны).
% Перенос строки ВНУТРИ \url{...} запрещён: пробел становится
% неразрывной частью слова, и URL уезжает за правый край.
\newcommand{\makemanual}{%
https://www.gnu.org/software/make/manual/%
make.html}

% Компактный титульный слайд.
% \insertsubtitle не используется: в шаблоне подзаголовок
% зафиксирован для другой темы презентации.
\setbeamertemplate{title page}{%
  \vfil
  \begin{center}
    {\large\inserttitle\par}%
    \vspace{0.45em}%
    {\small\textit{Сборка C/C++-проектов с помощью GNU Make}\par}%
    \vspace{0.4em}%
    {\small Курс «Системы сборки», Школа Разработчиков
      ООО «Прософт-Системы»\par}%
    \vspace{0.4em}%
    {\small Все примеры лекции — один проект из двух исходных
      файлов (\texttt{calc/})\par}%
    \vspace{1.0em}%
    {\small\insertauthor\par}%
    \vspace{0.25em}%
    {\small\insertinstitute\par}%
    \vspace{0.25em}%
    {\small\insertdate\par}%
  \end{center}%
  \vfil
}

\begin{document}

%------------------------------------------------------------------
% Слайд 1: титульный
%------------------------------------------------------------------
\begin{frame}
  \titlepage
\end{frame}

%------------------------------------------------------------------
% Слайд 2: зачем нужны системы сборки
%------------------------------------------------------------------
\begin{frame}{Зачем нужны системы сборки?}
  \small
  \begin{itemize}
    \item C/C++-программа собирается в два шага: компиляция
      каждого \texttt{.c} в \texttt{.o} и линковка
      \texttt{.o}-файлов в исполняемый файл.
    \item На маленьком проекте так делают руками, но легко
      ошибиться: неверный флаг, порядок файлов, забытая
      зависимость от заголовка.
    \item С ростом числа файлов (десятки, сотни) повторять
      команды руками нереально.
    \item Требования к сборке: воспроизводимость (одни и те же
      команды → один и тот же результат), скорость
      (пересобирать только изменившееся), единый источник
      правды о правилах сборки.
    \item Идея системы сборки — декларативность: описываем,
      \textbf{что} зависит от \textbf{чего}, а не пошаговый
      сценарий.
  \end{itemize}
\end{frame}

%------------------------------------------------------------------
% Слайд 3: рабочий пример — проект calc (часть 1: дерево)
%------------------------------------------------------------------
\begin{frame}[fragile]{Рабочий пример: проект \texttt{calc}}
  \small
  \myinputlisting[minted language=text]
    {presentations/01-gnu-make/src/}
    {project-tree.txt}
\end{frame}

%------------------------------------------------------------------
% Слайд 4: рабочий пример — проект calc (часть 2: calc.h)
%------------------------------------------------------------------
\begin{frame}[fragile]{Рабочий пример: проект \texttt{calc}}
  \small
  \myinputlisting[minted language=c]
    {presentations/01-gnu-make/src/}
    {calc.h}
  \begin{itemize}
    \item \texttt{calc.h} — единый заголовок: любой файл,
      включающий его, видит одни и те же объявления;
    \item \texttt{calc.c} и \texttt{main.c} включают
      \texttt{calc.h} — значит, при изменении заголовка
      пересобирать нужно \textbf{оба};
    \item артефакт сборки — исполняемый файл \texttt{calc};
      все Makefile лекции собирают именно его.
  \end{itemize}
\end{frame}

%------------------------------------------------------------------
% Слайд 5: наивное решение — shell-скрипт (часть 1)
%------------------------------------------------------------------
\begin{frame}[fragile]{Наивное решение: shell-скрипт}
  \small
  \myinputlisting[minted language=bash]
    {presentations/01-gnu-make/src/}
    {build.sh}
\end{frame}

%------------------------------------------------------------------
% Слайд 6: наивное решение — shell-скрипт (часть 2)
%------------------------------------------------------------------
\begin{frame}[fragile]{Наивное решение: shell-скрипт}
  \small
  \begin{itemize}
    \item полная пересборка при каждом запуске: компилируются
      оба \texttt{.c}-файла, даже если ничего не менялось;
    \item зависимости не отслеживаются — скрипт просто
      повторяет команды;
    \item флаги и компилятор захардкожены в тексте скрипта;
    \item не масштабируется: 50 исходников → 50 строк скрипта
      + 1 строка линковки;
    \item скрипт не отвечает на вопрос «что из проекта нужно
      пересобрать?».
  \end{itemize}
  \begin{terminalwindow}
$ ./build.sh
$ ./calc
2 + 3 = 5
10 - 4 = 6
6 * 7 = 42
20 / 5 = 4
  \end{terminalwindow}
\end{frame}

%------------------------------------------------------------------
% Слайд 7: правило GNU Make — анатомия
% Схема в две строки с буквальным маркером <TAB>: реального
% файла с таким содержимым не существует (табуляция в
% Makefile — настоящий символ TAB), поэтому листинг
% допустимо оставить in-place.
%------------------------------------------------------------------
\begin{frame}[fragile]{Правило GNU Make: анатомия}
  \small
  \begin{myinplacelisting}[minted language=text,
    minted options={fontsize=\small,breaklines,breakanywhere,
    linenos,numbersep=2pt}]
цели: зависимости
<TAB>рецепт
  \end{myinplacelisting}
  \footnotemark[1]\footnotemark[2]
  \begin{itemize}
    \item \textbf{цели} — что строим (файлы: \texttt{calc},
      \texttt{calc.o}, …);
    \item \textbf{зависимости} — что нужно для цели (файлы);
    \item \textbf{рецепт} — shell-команда, пересоздающая цель;
    \item рецепт начинается с \textbf{табуляции}, а не с
      пробелов (частая ошибка новичка);
    \item make выполняет рецепт, если цель отсутствует или
      \textbf{старше} любой из зависимостей (сравниваются
      метки времени).
  \end{itemize}
  \footnotetext[1]{%
    § 2.1 «What a Rule Looks Like» —
    \url{\makemanual\#What-a-Rule-Looks-Like}}
  \footnotetext[2]{%
    § 4.2 «Rule Syntax» —
    \url{\makemanual\#Rule-Syntax-1}}
\end{frame}

%------------------------------------------------------------------
% Слайд 8: первый Makefile — все правила руками
%------------------------------------------------------------------
\begin{frame}[fragile]{Первый Makefile: все правила руками}
  \small
  \myinputlisting[minted language=make]
    {presentations/01-gnu-make/src/}
    {Makefile.01}
  \begin{itemize}\setlength{\itemsep}{0pt}\setlength{\parsep}{0pt}
    \item максимальная явность: каждое действие записано
      полностью;
    \item \texttt{calc.h} в зависимостях \textbf{обоих}
      \texttt{.o} — правило «заголовок изменился → пересобрать
      всё, что его включает»;
    \item \texttt{make} без аргументов собирает первую цель
      файла — \texttt{calc}\footnotemark[3]\footnotemark[4].
  \end{itemize}
  \footnotetext[3]{%
    гл. 4 «Writing Rules» —
    \url{\makemanual\#Writing-Rules}}
  \footnotetext[4]{%
    § 2.2 «A Simple Makefile» —
    \url{\makemanual\#A-Simple-Makefile}}
\end{frame}

%------------------------------------------------------------------
% Слайд 9: инкрементальная сборка — что это даёт (часть 1)
%------------------------------------------------------------------
\begin{frame}[fragile]{Инкрементальная сборка: что это даёт}
  \small
  \begin{terminalwindow}
$ make
cc -Wall -O2 -c main.c -o main.o
cc -Wall -O2 -c calc.c -o calc.o
cc main.o calc.o -o calc
$ make
make: 'calc' is up to date.
$ touch calc.c
$ make
cc -Wall -O2 -c calc.c -o calc.o
cc main.o calc.o -o calc
  \end{terminalwindow}
  \footnotemark[5]
  \footnotetext[5]{%
    § 2.3 «How make Processes a Makefile» —
    \url{\makemanual\#How-make-Processes-a-Makefile}}
\end{frame}

%------------------------------------------------------------------
% Слайд 10: инкрементальная сборка — что это даёт (часть 2)
%------------------------------------------------------------------
\begin{frame}{Инкрементальная сборка: что это даёт}
  \small
  \begin{itemize}
    \item повторный запуск: make сравнивает метки времени и не
      выполняет ни одной команды;
    \item изменился только \texttt{calc.c} → перекомпилирован
      только он; \texttt{main.o} не тронут;
    \item линковка всё же выполнена: \texttt{calc} оказался
      старше нового \texttt{calc.o};
    \item вывод \texttt{make} — живая «отчётность»: видно, что
      именно делалось.
  \end{itemize}
\end{frame}

%------------------------------------------------------------------
% Слайд 11: автоматические переменные (часть 1)
%------------------------------------------------------------------
\begin{frame}[fragile]{Упрощение 1: автоматические переменные}
  \small
  \myinputlisting[minted language=make]
    {presentations/01-gnu-make/src/}
    {Makefile.02}
  \footnotemark[6]
  \footnotetext[6]{%
    § 10.5.3 «Automatic Variables» —
    \url{\makemanual\#Automatic-Variables}}
\end{frame}

%------------------------------------------------------------------
% Слайд 12: автоматические переменные (часть 2)
%------------------------------------------------------------------
\begin{frame}{Упрощение 1: автоматические переменные}
  \small
  \begin{center}
    \small
    \begin{tabular}{@{}p{2.0cm}p{3.4cm}p{6.8cm}@{}}
      \textbf{Переменная} & \textbf{Что подставляет}
        & \textbf{Пример из нашего Makefile} \\
      \texttt{\$@} & имя цели & в правиле \texttt{main.o: …}
        это \texttt{main.o} \\
      \texttt{\$<} & \textbf{первая} зависимость & исходный
        \texttt{.c}-файл \\
      \texttt{\$\^} & \textbf{все} зависимости & в правиле
        линковки — оба \texttt{.o} \\
    \end{tabular}
  \end{center}
  \begin{itemize}
    \item \small правило перестаёт «помнить» имена файлов:
      проще переименовать, легче копировать в новый проект.
  \end{itemize}
\end{frame}

%------------------------------------------------------------------
% Слайд 13: шаблонные правила
%------------------------------------------------------------------
\begin{frame}[fragile]{Упрощение 2: шаблонные правила}
  \small
  \myinputlisting[minted language=make]
    {presentations/01-gnu-make/src/}
    {Makefile.03}
  \footnotemark[7]\footnotemark[8]
  \begin{itemize}
    \item \texttt{\%} — «ствол» имени: для цели
      \texttt{main.o} ствол — \texttt{main}, поэтому
      \texttt{\$<} = \texttt{main.c};
    \item \texttt{calc.h} оставлен в зависимостях шаблона:
      make не знает, что наши исходники включают этот
      заголовок;
    \item новый исходник: положить \texttt{.c}-файл и добавить
      его имя в правило линковки — отдельное правило на него
      не пишется.
  \end{itemize}
  \footnotetext[7]{%
    гл. 10 «Using Implicit Rules» —
    \url{\makemanual\#Using-Implicit-Rules}}
  \footnotetext[8]{%
    § 10.5.1 «Introduction to Pattern Rules» —
    \url{\makemanual\#Introduction-to-Pattern-Rules}}
\end{frame}

%------------------------------------------------------------------
% Слайд 14: встроенные неявные правила (часть 1)
%------------------------------------------------------------------
\begin{frame}[fragile]{Встроенные неявные правила: что умеет
  make «из коробки»}
  \small
  \myinputlisting[minted language=make]
    {presentations/01-gnu-make/src/}
    {Makefile.03b}
  \begin{terminalwindow}
$ make
cc    -c -o calc.o calc.c
cc    -c -o main.o main.c
cc   calc.o main.o   -o calc
$ touch calc.h
$ make
make: 'calc' is up to date.
  \end{terminalwindow}
  \footnotemark[9]
  \footnotetext[9]{%
    § 10.2 «Catalogue of Built-In Rules» —
    \url{\makemanual\#Catalogue-of-Built_002dIn-Rules}}
\end{frame}

%------------------------------------------------------------------
% Слайд 15: встроенные неявные правила (часть 2)
%------------------------------------------------------------------
\begin{frame}{Встроенные неявные правила: что умеет
  make «из коробки»}
  \small
  \begin{itemize}
    \item make знает встроенное правило \texttt{\%.o: \%.c}
      (и правило линковки) — сборка «прошла» вообще без
      правил;
    \item но флаги \texttt{-Wall -O2} исчезли, а
      \texttt{calc.h} не входит в зависимости;
    \item изменение заголовка make не заметил — сборка стала
      \textbf{неправильной} (бинарник устарел, но make этого
      не знает);
    \item вывод: встроенные правила — удобная отправная
      точка; для серьёзного проекта пишем свои шаблоны.
  \end{itemize}
\end{frame}

%------------------------------------------------------------------
% Слайд 16: финальный Makefile простого проекта (часть 1)
%------------------------------------------------------------------
\begin{frame}[fragile]{Финальный Makefile простого проекта}
  \small
  \myinputlisting[minted language=make]
    {presentations/01-gnu-make/src/}
    {Makefile.04}
  \footnotemark[10]
  \footnotetext[10]{%
    гл. 6 «How to Use Variables» —
    \url{\makemanual\#How-to-Use-Variables}}
\end{frame}

%------------------------------------------------------------------
% Слайд 17: финальный Makefile простого проекта (часть 2)
%------------------------------------------------------------------
\begin{frame}[fragile]{Финальный Makefile простого проекта}
  \small
  \begin{columns}[T]
    \begin{column}{0.55\textwidth}
      \begin{terminalwindow}
$ touch calc.h
$ make
cc -Wall -O2 -c main.c -o main.o
cc -Wall -O2 -c calc.c -o calc.o
cc -Wall -O2 main.o calc.o -o calc
      \end{terminalwindow}
    \end{column}
    \begin{column}{0.40\textwidth}
      \begin{itemize}
        \item компилятор и флаги вынесены в переменные —
          меняются в одном месте;
        \item изменение \texttt{calc.h} → перекомпилированы
          \textbf{оба} \texttt{.o}: слежение за заголовком
          работает;
        \item итог преобразования: 3 раздутых правила → 1
          шаблон + 2 переменные;
        \item 8 строк — вся сборочная система маленького
          проекта.
      \end{itemize}
    \end{column}
  \end{columns}
\end{frame}

%------------------------------------------------------------------
% Слайд 18: служебные цели all, clean, help (часть 1: листинг,
% строки 1–10 из 18)
% 18-строчный листинг с заголовком кадра и подписью файла не
% влезает в один кадр, поэтому разбит на два: строки 1–10.
%------------------------------------------------------------------
\begin{frame}[fragile]{Служебные цели: \texttt{all},
  \texttt{clean}, \texttt{help}}
  \small
  \myinputlisting[minted language=make,
    minted options={fontsize=\small,breaklines,breakanywhere,
    linenos,numbersep=2pt,lastline=10}]
    {presentations/01-gnu-make/src/}
    {Makefile.05}
\end{frame}

%------------------------------------------------------------------
% Слайд 19: служебные цели all, clean, help (часть 2: листинг,
% строки 11–18)
%------------------------------------------------------------------
\begin{frame}[fragile]{Служебные цели: \texttt{all},
  \texttt{clean}, \texttt{help}}
  \small
  \myinputlisting[minted language=make,
    minted options={fontsize=\small,breaklines,breakanywhere,
    linenos,numbersep=2pt,firstline=11}]
    {presentations/01-gnu-make/src/}
    {Makefile.05}
\end{frame}

%------------------------------------------------------------------
% Слайд 20: служебные цели all, clean, help (часть 3:
% пояснения)
%------------------------------------------------------------------
\begin{frame}{Служебные цели: \texttt{all}, \texttt{clean},
  \texttt{help}}
  \small
  \footnotemark[11]\footnotemark[12]
  \begin{itemize}\setlength{\itemsep}{0pt}\setlength{\parsep}{0pt}
    \item \texttt{all} — первая цель файла; \texttt{make} без
      аргументов = «построй всё»;
    \item \texttt{clean} — удаляет артефакты
      (\texttt{rm -f …}), возвращает каталог в исходное
      состояние;
    \item \texttt{help} — печатает список целей; \texttt{@} в
      начале команды подавляет её вывод\footnotemark[13];
    \item строка \texttt{.PHONY: all clean help} — «это не
      файлы» (почему — на следующем слайде).
  \end{itemize}
  \footnotetext[11]{%
    § 2.7 «Rules for Cleaning the Directory» —
    \url{\makemanual\#Rules-for-Cleaning-the-Directory}}
  \footnotetext[12]{%
    § 16.6 «Standard Targets for Users» —
    \url{\makemanual\#Standard-Targets-for-Users}}
  \footnotetext[13]{%
    § 5.2 «Recipe Echoing» —
    \url{\makemanual\#Recipe-Echoing}
    (о символе \texttt{@})}
\end{frame}

%------------------------------------------------------------------
% Слайд 21: зачем нужен .PHONY (часть 1)
%------------------------------------------------------------------
\begin{frame}[fragile]{Зачем нужен \texttt{.PHONY}}
  \small
  \begin{itemize}
    \item Сценарий: в каталоге появляется файл с именем
      \texttt{clean} (ошибка имени, временный файл, …).
    \item Демонстрация на \texttt{Makefile.05bug} — тот же
      Makefile \textbf{без} строки \texttt{.PHONY}.
  \end{itemize}
  \begin{terminalwindow}
$ touch clean
$ make clean
make: 'clean' is up to date.
$ ls
build.sh  calc  calc.c  calc.h  calc.o  clean  main.c  main.o  …
  \end{terminalwindow}
\end{frame}

%------------------------------------------------------------------
% Слайд 22: зачем нужен .PHONY (часть 2)
%------------------------------------------------------------------
\begin{frame}{Зачем нужен \texttt{.PHONY}}
  \small
  \begin{itemize}
    \item цель \texttt{clean} совпала по имени с файлом → make
      рассматривает её как файл;
    \item файл существует и зависимостей не имеет → «в
      актуальном состоянии», рецепт \textbf{не выполняется} —
      артефакты не удалены;
    \item \texttt{.PHONY} объявляет цели фиктивными: файл с
      таким именем игнорируется, рецепт выполняется всегда;
    \item правило: все «псевдоцели» (\texttt{all},
      \texttt{clean}, …) должны входить в
      \texttt{.PHONY}\footnotemark[14]\footnotemark[15].
  \end{itemize}
  \footnotetext[14]{%
    § 4.6 «Phony Targets» —
    \url{\makemanual\#Phony-Targets-1}}
  \footnotetext[15]{%
    § 4.9 «Special Built-in Target Names» —
    \url{\makemanual\#Special-Built_002din-Target-Names}}
\end{frame}

%------------------------------------------------------------------
% Слайд 23: статическая библиотека libcalc.a (часть 1)
% Заголовок уменьшен до \normalsize: длинный заголовок в \large
% перекрывает логотип в правом верхнем углу.
%------------------------------------------------------------------
\setbeamerfont{frametitle}{size=\normalsize}
\begin{frame}{Проект посложнее: статическая библиотека
  \texttt{libcalc.a}}
  \small
  \begin{itemize}
    \item Дерево не меняется, но роли: \texttt{calc.c} — код
      \textbf{библиотеки}, \texttt{main.c} — \textbf{клиент};
    \item артефакты: \texttt{libcalc.a} (статическая
      библиотека) + \texttt{calc} (клиентский exe).
  \end{itemize}
\end{frame}

%------------------------------------------------------------------
% Слайд 24: статическая библиотека libcalc.a (часть 2:
% листинг, строки 1–10 из 18)
% 18-строчный листинг с заголовком кадра и подписью файла не
% влезает в один кадр, поэтому разбит на два: строки 1–10.
%------------------------------------------------------------------
\begin{frame}[fragile]{Проект посложнее: статическая
  библиотека \texttt{libcalc.a}}
  \small
  \myinputlisting[minted language=make,
    minted options={fontsize=\small,breaklines,breakanywhere,
    linenos,numbersep=2pt,lastline=10}]
    {presentations/01-gnu-make/src/}
    {Makefile.06}
\end{frame}

%------------------------------------------------------------------
% Слайд 25: статическая библиотека libcalc.a (часть 3:
% листинг, строки 11–18)
%------------------------------------------------------------------
\begin{frame}[fragile]{Проект посложнее: статическая
  библиотека \texttt{libcalc.a}}
  \small
  \myinputlisting[minted language=make,
    minted options={fontsize=\small,breaklines,breakanywhere,
    linenos,numbersep=2pt,firstline=11}]
    {presentations/01-gnu-make/src/}
    {Makefile.06}
\end{frame}

%------------------------------------------------------------------
% Слайд 26: статическая библиотека libcalc.a (часть 4)
%------------------------------------------------------------------
\begin{frame}[fragile]{Проект посложнее: статическая
  библиотека \texttt{libcalc.a}}
  \small
  \begin{columns}[T]
    \begin{column}{0.55\textwidth}
      \begin{terminalwindow}
$ make
cc -Wall -O2 -c main.c -o main.o
cc -Wall -O2 -c calc.c -o calc.o
ar rcs libcalc.a calc.o
cc main.o libcalc.a -o calc
$ touch calc.c
$ make
cc -Wall -O2 -c calc.c -o calc.o
ar rcs libcalc.a calc.o
cc main.o libcalc.a -o calc
      \end{terminalwindow}
    \end{column}
    \begin{column}{0.40\textwidth}
      \begin{itemize}
        \item \texttt{ar rcs} — положить \texttt{calc.o} в
          архив и построить индекс символов;
        \item изменился \texttt{calc.c} → пересобрана
          библиотека и перелинкован клиент, \texttt{main.o} не
          тронут;
        \item шаблонное правило \texttt{\%.o: \%.c calc.h}
          одинаково работает для кода библиотеки и клиента;
        \item статическая библиотека «вшивается» в
          исполняемый файл при линковке.
      \end{itemize}
    \end{column}
  \end{columns}
\end{frame}
\setbeamerfont{frametitle}{size=\Large}

%------------------------------------------------------------------
% Слайд 27: динамическая библиотека libcalc.so (часть 1:
% листинг, строки 1–10 из 18)
% 18-строчный листинг с заголовком кадра и подписью файла не
% влезает в один кадр, поэтому разбит на два: строки 1–10.
%------------------------------------------------------------------
\begin{frame}[fragile]{Динамическая библиотека
  \texttt{libcalc.so}}
  \small
  \myinputlisting[minted language=make,
    minted options={fontsize=\small,breaklines,breakanywhere,
    linenos,numbersep=2pt,lastline=10}]
    {presentations/01-gnu-make/src/}
    {Makefile.07}
\end{frame}

%------------------------------------------------------------------
% Слайд 28: динамическая библиотека libcalc.so (часть 2:
% листинг, строки 11–18)
%------------------------------------------------------------------
\begin{frame}[fragile]{Динамическая библиотека
  \texttt{libcalc.so}}
  \small
  \myinputlisting[minted language=make,
    minted options={fontsize=\small,breaklines,breakanywhere,
    linenos,numbersep=2pt,firstline=11}]
    {presentations/01-gnu-make/src/}
    {Makefile.07}
\end{frame}

%------------------------------------------------------------------
% Слайд 29: динамическая библиотека libcalc.so (часть 3:
% терминал)
%------------------------------------------------------------------
\begin{frame}[fragile]{Динамическая библиотека
  \texttt{libcalc.so}}
  \small
  \begin{terminalwindow}
$ make
cc -Wall -O2 -fPIC -c main.c -o main.o
cc -Wall -O2 -fPIC -c calc.c -o calc.o
cc -shared calc.o -o libcalc.so
cc main.o libcalc.so -o calc
$ ./calc
./calc: error while loading shared libraries:
    libcalc.so: cannot open shared object file: No such file or directory
$ LD_LIBRARY_PATH=. ./calc
2 + 3 = 5
10 - 4 = 6
6 * 7 = 42
20 / 5 = 4
  \end{terminalwindow}
\end{frame}

%------------------------------------------------------------------
% Слайд 30: динамическая библиотека libcalc.so (часть 4)
%------------------------------------------------------------------
\begin{frame}{Динамическая библиотека \texttt{libcalc.so}}
  \small
  \begin{itemize}
    \item код библиотеки компилируется с \texttt{-fPIC}
      (независимый от адреса), сама библиотека линкуется с
      \texttt{-shared};
    \item исполняемый файл хранит лишь \textbf{имя}
      \texttt{libcalc.so} — загрузка происходит в момент
      запуска;
    \item загрузчику нужно знать, где искать:
      \texttt{LD\_LIBRARY\_PATH=.} (в production — rpath или
      системные каталоги);
    \item отличие от статической: библиотека не вшита в
      \texttt{calc}, общий код обновляется заменой
      \texttt{.so}.
  \end{itemize}
\end{frame}

%------------------------------------------------------------------
% Слайд 31: встроенные функции (часть 1: листинг,
% строки 1–14 из 22)
% Исходный слайд переполнял кадр, листинг разбит; строки 1–14
% идут в кадр с заголовком.
%------------------------------------------------------------------
\begin{frame}[fragile]{Встроенные функции: \texttt{\$(wildcard)},
  \texttt{\$(patsubst)}, \texttt{\$(shell)}}
  \small
  \myinputlisting[minted language=make,
    minted options={fontsize=\small,breaklines,breakanywhere,
    linenos,numbersep=2pt,lastline=14}]
    {presentations/01-gnu-make/src/}
    {Makefile.08}
\end{frame}

%------------------------------------------------------------------
% Слайд 32: встроенные функции (часть 2: листинг,
% строки 15–22, терминал)
%------------------------------------------------------------------
\begin{frame}[fragile]{Встроенные функции: \texttt{\$(wildcard)},
  \texttt{\$(patsubst)}, \texttt{\$(shell)}}
  \small
  \myinputlisting[minted language=make,
    minted options={fontsize=\small,breaklines,breakanywhere,
    linenos,numbersep=2pt,firstline=15}]
    {presentations/01-gnu-make/src/}
    {Makefile.08}
  \begin{terminalwindow}
$ make info
Compiler: clang version 21.1.8
Sources : calc.c main.c
  \end{terminalwindow}
\end{frame}

%------------------------------------------------------------------
% Слайд 33: встроенные функции (часть 3: таблица, строки 1–2,
% сноски 16–18)
% Вся таблица + все 4 сноски не влезает ни в один кадр
% (проверено измерением), таблица разбита: строки 1–2 идут в
% кадр с заголовком; сноска 19 — в кадре 34, у строки про
% \$(shell). Отметки сносок 16–18 стоят в столбце «Функция»
% на строках своих функций: отдельная строка с отметками
% переполняла кадр на 12,3pt (высота строки \small).
% \topsep=0pt убирает клей \topsep (10pt сверху/снизу)
% вокруг center'а — без этого не хватало 2,8pt по высоте.
%------------------------------------------------------------------
\begin{frame}{Встроенные функции: \texttt{\$(wildcard)},
  \texttt{\$(patsubst)}, \texttt{\$(shell)}}
  \small
  \setlength{\topsep}{0pt}
  \setlength{\tabcolsep}{5pt}
  \begin{center}
    \begin{tabular}{@{}p{5.0cm}p{3.8cm}p{4.5cm}@{}}
      \textbf{Функция} & \textbf{Что делает}
        & \textbf{Где используется в \texttt{calc}} \\
      \texttt{\$(wildcard *.c)}\footnotemark[16]\footnotemark[17]
      & shell-подобный глоб: список
        всех \texttt{.c} в каталоге & список исходников
        ведётся автоматически, новый файл «появляется» в
        сборке сам \\
      \texttt{\$(patsubst \%.c,\%.o,\$(SRCS))}\footnotemark[18]
        & подстановка по образцу: \texttt{main.c → main.o}
        & из списка исходников автоматически получается список
        объектов \\
    \end{tabular}
  \end{center}
  \footnotetext[16]{%
    гл. 8 «Functions for Transforming Text» —
    \url{\makemanual\#Functions-for-Transforming-Text}}
  \footnotetext[17]{%
    § 8.3 «Functions for File Names» —
    \url{\makemanual\#Functions-for-File-Names}
    (\texttt{\$(wildcard)})}
  \footnotetext[18]{%
    § 8.2 «Functions for String Substitution and Analysis» —
    \url{\makemanual\#Functions-for-String-Substitution-and-Analysis}
    (\texttt{\$(patsubst)})}
\end{frame}

%------------------------------------------------------------------
% Слайд 34: встроенные функции (часть 4: строка 3 таблицы,
% сноска 19)
%------------------------------------------------------------------
\begin{frame}{Встроенные функции: \texttt{\$(wildcard)},
  \texttt{\$(patsubst)}, \texttt{\$(shell)}}
  \small
  \setlength{\tabcolsep}{5pt}
  \begin{center}
    \begin{tabular}{@{}p{5.0cm}p{3.8cm}p{4.5cm}@{}}
      \textbf{Функция} & \textbf{Что делает}
        & \textbf{Где используется в \texttt{calc}} \\
      \texttt{\$(shell \$(CC) --version \textbar head -n1)}\footnotemark[19]
        & shell-команда, выполняемая при разборе Makefile
        & версия компилятора в цели \texttt{info} \\
    \end{tabular}
  \end{center}
  \begin{itemize}
    \item результат функций вычисляется при раскрытии
      переменной.
  \end{itemize}
  \footnotetext[19]{%
    § 8.14 «The shell Function» —
    \url{\makemanual\#The-shell-Function}}
\end{frame}

%------------------------------------------------------------------
% Слайд 35: пользовательские функции (часть 1: листинг,
% строки 1–13 из 23)
% 23-строчный листинг с заголовком кадра не влезает в один
% кадр, поэтому разбит на два: строки 1–13. Сноски 20–22
% стоят в кадре 36, в конце второй части.
%------------------------------------------------------------------
\begin{frame}[fragile]{Пользовательские функции: \texttt{define},
  \texttt{call}, \texttt{eval}}
  \small
  \myinputlisting[minted language=make,
    minted options={fontsize=\small,breaklines,breakanywhere,
    linenos,numbersep=2pt,lastline=13}]
    {presentations/01-gnu-make/src/}
    {Makefile.09}
\end{frame}

%------------------------------------------------------------------
% Слайд 36: пользовательские функции (часть 2: листинг,
% строки 14–23)
%------------------------------------------------------------------
\begin{frame}[fragile]{Пользовательские функции: \texttt{define},
  \texttt{call}, \texttt{eval}}
  \small
  \myinputlisting[minted language=make,
    minted options={fontsize=\small,breaklines,breakanywhere,
    linenos,numbersep=2pt,firstline=14}]
    {presentations/01-gnu-make/src/}
    {Makefile.09}
  \footnotemark[20]\footnotemark[21]\footnotemark[22]
  \footnotetext[20]{%
    § 8.1 «Function Call Syntax» —
    \url{\makemanual\#Function-Call-Syntax}}
  \footnotetext[21]{%
    § 8.8 «The call Function» —
    \url{\makemanual\#The-call-Function}}
  \footnotetext[22]{%
    § 8.10 «The eval Function» —
    \url{\makemanual\#The-eval-Function}}
\end{frame}

%------------------------------------------------------------------
% Слайд 37: пользовательские функции (часть 3: пояснения)
%------------------------------------------------------------------
\begin{frame}{Пользовательские функции: \texttt{define},
  \texttt{call}, \texttt{eval}}
  \small
  \begin{itemize}
    \item \texttt{define … endef} — многострочная переменная
      (тело функции) с плейсхолдерами \texttt{\$(1)},
      \texttt{\$(2)};
    \item \texttt{\$(call BUILD\_STATIC\_LIB,libcalc,calc.o)}
      — подстановка аргументов в тело;
    \item \texttt{\$(eval …)} — результат разбирается как код
      Makefile: правило \texttt{libcalc.a: calc.o}
      \textbf{генерируется} при разборе файла;
    \item \texttt{\$\$@} / \texttt{\$\$\^} в теле —
      экранированный доллар: в рецепт попадут живые
      \texttt{\$@} и \texttt{\$\^};
    \item одной функцией можно сгенерировать правила для
      сколь угодно многих библиотек (вызвать N раз);
    \item тонкость: \texttt{\$(eval …)} в начале файла делает
      сгенерированную цель \textbf{дефолтной} (первая цель
      файла) — поэтому \texttt{all: calc} стоит выше вызова.
  \end{itemize}
\end{frame}

%------------------------------------------------------------------
% Слайд 38: итоги
%------------------------------------------------------------------
\begin{frame}{Итоги}
  \small
  \begin{itemize}
    \item Правило = \texttt{цели: зависимости} + рецепт
      (TAB); make принимает решения по меткам времени — отсюда
      инкрементальная сборка.
    \item Автоматические переменные (\texttt{\$@},
      \texttt{\$<}, \texttt{\$\^}) и шаблонные правила убирают
      шаблонный код; встроенные правила — отправная точка, а
      не замена своим.
    \item Служебные цели \texttt{all}, \texttt{clean},
      \texttt{help} + \texttt{.PHONY} защищают от коллизий с
      именами файлов.
    \item Функции (\texttt{\$(wildcard)},
      \texttt{\$(patsubst)}, \texttt{\$(shell)},
      \texttt{define/call/eval}) масштабируют сборку на
      большие проекты.
    \item Руководство GNU Make: \url{\makemanual}\footnotemark[23]
  \end{itemize}
  \footnotetext[23]{%
    Руководство GNU Make, гл. 2 «Makefiles» —
    \url{\makemanual\#Makefiles}}
\end{frame}

\end{document}
